% GENERATED FILE. Edit canonical lessons, then run npm run generate:books.
% canonical-source-hash: fnv1a64:d4809a5cef61a583
% canonical-lessons: ML-C201-atupole, ML-C201-olike, ML-C201-bhagyavasal, ML-C201-nirbhagyavasal, ML-C201-ennaneyenkilum

\chapter{Likewise, Except, Fortunately, Unfortunately, Somehow}
\label{ch:ml-likewise-except-fortunately-unfortunately-somehow}
\hlchaptermodality{201}

\begin{chapteropening}
\textbf{By the end of this chapter:} \emph{I can say likewise, except, fortunately, unfortunately and somehow.}

Complete the last lesson of chapter 201: I can say likewise, except, fortunately, unfortunately and somehow.
\end{chapteropening}

\section[\texorpdfstring{\ml{അതുപോലെ} (atupōle)}{അതുപോലെ (atupōle)}]{\ml{അതുപോലെ} (atupōle) — likewise, in the same way}

\label{lesson:ML-C201-atupole}



\begin{quote}
Before the new one: say the Malayalam for a solution, then the Malayalam for according to.
\end{quote}

\subsection*{You'll want to know: \ml{അതുപോലെ}}
\textbf{\ml{അതുപോലെ}} — \emph{atupōle} — \textquotedblleft{}likewise, in the same way\textquotedblright{}.

\begin{grammarlens}[title={What you've built}]
One more word for this chapter.
\end{grammarlens}

\subsection*{Your turn}
\begin{itemize}
\raggedright
  \item \emph{Say it:} \emph{atupōle}
  \item \emph{Say it:} \emph{atupōle}, once more
  \item \emph{Say it:} say \emph{atupōle}
  \item \emph{Recall:} say the Malayalam for a solution, then the Malayalam for according to, then say \emph{atupōle} again
\end{itemize}

\begin{tcolorbox}[breakable,colback=teal!4,colframe=teal!35!black,title={Before you move on}]
What does \ml{അതുപോലെ} mean? (\textquotedblleft{}Likewise, in the same way\textquotedblright{}.) Say it once more.
\end{tcolorbox}
\section[\texorpdfstring{\ml{ഒഴികെ} (oḻike)}{ഒഴികെ (oḻike)}]{\ml{ഒഴികെ} (oḻike) — except}

\label{lesson:ML-C201-olike}



\begin{quote}
Before the new one: say the Malayalam for according to, then the Malayalam for likewise.
\end{quote}

\subsection*{You'll want to know: \ml{ഒഴികെ}}
\textbf{\ml{ഒഴികെ}} — \emph{oḻike} — \textquotedblleft{}except\textquotedblright{}.

\begin{grammarlens}[title={What you've built}]
One more word for this chapter.
\end{grammarlens}

\subsection*{Your turn}
\begin{itemize}
\raggedright
  \item \emph{Say it:} \emph{oḻike}
  \item \emph{Say it:} \emph{oḻike}, once more
  \item \emph{Say it:} say \emph{oḻike}
  \item \emph{Recall:} say the Malayalam for according to, then the Malayalam for likewise, then say \emph{oḻike} again
\end{itemize}

\begin{tcolorbox}[breakable,colback=teal!4,colframe=teal!35!black,title={Before you move on}]
What does \ml{ഒഴികെ} mean? (\textquotedblleft{}Except\textquotedblright{}.) Say it once more.
\end{tcolorbox}
\section[\texorpdfstring{\ml{ഭാഗ്യവശാൽ} (bhāgyavaśāl)}{ഭാഗ്യവശാൽ (bhāgyavaśāl)}]{\ml{ഭാഗ്യവശാൽ} (bhāgyavaśāl) — fortunately}

\label{lesson:ML-C201-bhagyavasal}



\begin{quote}
Before the new one: say the Malayalam for likewise, then the Malayalam for except.
\end{quote}

\subsection*{You'll want to know: \ml{ഭാഗ്യവശാൽ}}
\textbf{\ml{ഭാഗ്യവശാൽ}} — \emph{bhāgyavaśāl} — \textquotedblleft{}fortunately\textquotedblright{}.

\begin{grammarlens}[title={What you've built}]
One more word for this chapter.
\end{grammarlens}

\subsection*{Your turn}
\begin{itemize}
\raggedright
  \item \emph{Say it:} \emph{bhāgyavaśāl}
  \item \emph{Say it:} \emph{bhāgyavaśāl}, once more
  \item \emph{Say it:} say \emph{bhāgyavaśāl}
  \item \emph{Recall:} say the Malayalam for likewise, then the Malayalam for except, then say \emph{bhāgyavaśāl} again
\end{itemize}

\begin{tcolorbox}[breakable,colback=teal!4,colframe=teal!35!black,title={Before you move on}]
What does \ml{ഭാഗ്യവശാൽ} mean? (\textquotedblleft{}Fortunately\textquotedblright{}.) Say it once more.
\end{tcolorbox}
\section[\texorpdfstring{\ml{നിർഭാഗ്യവശാൽ} (nirbhāgyavaśāl)}{നിർഭാഗ്യവശാൽ (nirbhāgyavaśāl)}]{\ml{നിർഭാഗ്യവശാൽ} (nirbhāgyavaśāl) — unfortunately}

\label{lesson:ML-C201-nirbhagyavasal}



\begin{quote}
Before the new one: say the Malayalam for except, then the Malayalam for fortunately.
\end{quote}

\subsection*{You'll want to know: \ml{നിർഭാഗ്യവശാൽ}}
\textbf{\ml{നിർഭാഗ്യവശാൽ}} — \emph{nirbhāgyavaśāl} — \textquotedblleft{}unfortunately\textquotedblright{}.

\begin{grammarlens}[title={What you've built}]
One more word for this chapter.
\end{grammarlens}

\subsection*{Your turn}
\begin{itemize}
\raggedright
  \item \emph{Say it:} \emph{nirbhāgyavaśāl}
  \item \emph{Say it:} \emph{nirbhāgyavaśāl}, once more
  \item \emph{Say it:} say \emph{nirbhāgyavaśāl}
  \item \emph{Recall:} say the Malayalam for except, then the Malayalam for fortunately, then say \emph{nirbhāgyavaśāl} again
\end{itemize}

\begin{tcolorbox}[breakable,colback=teal!4,colframe=teal!35!black,title={Before you move on}]
What does \ml{നിർഭാഗ്യവശാൽ} mean? (\textquotedblleft{}Unfortunately\textquotedblright{}.) Say it once more.
\end{tcolorbox}
\section[\texorpdfstring{\ml{എങ്ങനെയെങ്കിലും} (eṅṅaneyeṅkiluṁ)}{എങ്ങനെയെങ്കിലും (eṅṅaneyeṅkiluṁ)}]{\ml{എങ്ങനെയെങ്കിലും} (eṅṅaneyeṅkiluṁ) — somehow}

\label{lesson:ML-C201-ennaneyenkilum}



\begin{quote}
Before the new one: say the Malayalam for fortunately, then the Malayalam for unfortunately.
\end{quote}

\subsection*{You'll want to know: \ml{എങ്ങനെയെങ്കിലും}}
\textbf{\ml{എങ്ങനെയെങ്കിലും}} — \emph{eṅṅaneyeṅkiluṁ} — \textquotedblleft{}somehow\textquotedblright{}.

\begin{grammarlens}[title={What you've built}]
One more word for this chapter.
\end{grammarlens}

\subsection*{Your turn}
\begin{itemize}
\raggedright
  \item \emph{Say it:} \emph{eṅṅaneyeṅkiluṁ}
  \item \emph{Say it:} \emph{eṅṅaneyeṅkiluṁ}, once more
  \item \emph{Say it:} say \emph{eṅṅaneyeṅkiluṁ}
  \item \emph{Recall:} say the Malayalam for fortunately, then the Malayalam for unfortunately, then say \emph{eṅṅaneyeṅkiluṁ} again
\end{itemize}

\begin{tcolorbox}[breakable,colback=teal!4,colframe=teal!35!black,title={Before you move on}]
What does \ml{എങ്ങനെയെങ്കിലും} mean? (\textquotedblleft{}Somehow\textquotedblright{}.) Say it once more.
\end{tcolorbox}
